\chapter{Other works \& achievements}\label{chapter:other}

This chapter opens with a brief overview of other papers I co-authored during my PhD studies. These works either fall outside the thematic scope of this thesis or are part of ongoing research and are therefore omitted from the main text. Nevertheless, as they are relevant to my broader research area, they are included here. The chapter then concludes with a summary of my other achievements, including complementary research activities and outreach initiatives aimed at advancing deep learning.

\section{Publications}

\begin{itemize}
    \item \otherpubentry{kargin2026sparrta}
    
    Summary:
    \begin{quote}
    Visual foundation models recognise semantic content in images, but their ability to represent spatial relations between objects remains less well understood. The study introduces the Spatial Relation Recognition Task (SpaRRTa), a benchmark that evaluates this ability independently of semantic classification and precise geometric prediction. A controllable synthetic data pipeline generates photorealistic scenes with varied object arrangements and spatial annotations. The evaluation compares visual foundation models using different probing methods to assess how spatial information can be extracted from their representations. Results show that spatial relations are primarily encoded in patch-level features and can be obscured by global pooling, causing standard linear evaluation to underestimate spatial awareness. Relations defined relative to another object remain challenging across model families. These findings establish SpaRRTa as a diagnostic benchmark for distinguishing limitations of visual representations from limitations of the methods used to read out spatial information.
    \end{quote}
    \item {\emergencystretch=1em \otherpubentry{sroka2025ai}\par}
    
    Summary:
    \begin{quote}
    Identifying pathogens from blood samples is an important step in sepsis diagnosis, but conventional microbiological procedures can be time-consuming. The study develops a deep learning pipeline for species-level identification from microscopic images of Gram-stained smears of positive blood samples. It combines Cellpose 3 for microbial cell segmentation with attention-based deep multiple instance learning for classification, allowing information from multiple cells to contribute to each prediction. The evaluation uses 16,637 images covering 14 bacterial and three yeast-like fungal species. The model achieves classification accuracies of 77.15\% for bacteria and 71.39\% for fungi. Errors are associated with morphological similarities between species and variation within individual species. The results support the potential of microscopy-based microbial identification as an aid to diagnosis, while indicating that further optimisation and more diverse training data are needed to improve classification reliability.
    \end{quote}
    \item \otherpubentry{pardyl2023complung}
    
    Summary:
    \begin{quote}
    Computer-aided analysis of lung computed tomography (CT) scans often addresses individual processing stages, leaving their integration into patient-level assessment unresolved. The study introduces CompLung, a pipeline combining lung segmentation, suspicious-region detection, patch representation learning, and patient-level classification. It uses multiple instance learning to aggregate information from candidate regions into a prediction for the complete scan. The evaluation uses LIDC-IDRI, extended with radiologist-produced lung segmentation masks made available for further research. For patient-level classification using a malignancy-score threshold greater than three, CompLung achieves an area under the receiver operating characteristic curve of 0.90, compared with 0.83 for the strongest reported baselines. Qualitative analyses also examine the regions contributing to predictions. The study demonstrates the benefit of integrating the processing stages for prediction of radiologist-derived labels and provides segmentation annotations to support subsequent work on automated CT analysis.
    \end{quote}
    \item \otherpubentry{struski2023propml}
    
    Summary:
    \begin{quote}
    Partial multi-label learning addresses training data in which each instance has a candidate label set containing both correct and incorrect labels. Existing approaches often require a separate procedure to identify the relevant labels, increasing training complexity. The study introduces ProPML, a probabilistic loss that adapts binary cross-entropy to this setting. It encourages predictions within the candidate set while penalising predictions outside it, and can be used with deep networks through a change to the training objective. The evaluation covers datasets with naturally ambiguous labels and artificially corrupted labels, including image classification on VOC2007 and COCO2014, across several noise levels. ProPML compares favourably with existing methods, with particularly strong results when candidate sets contain more incorrect labels. These findings support a simple loss-based approach to learning from ambiguous annotations without a separate label-disambiguation procedure.
    \end{quote}
    \item \otherpubentry{pardyl2022automating}

    Summary:
    \begin{quote}
    Patient-level diagnostics from lung computed tomography (CT) scans requires models that can operate with different amounts of annotation. The study introduces MilLung and AutoLung and compares them with the existing DeepLung baseline to examine the effect of supervision on automated assessment. DeepLung uses nodule locations and malignancy annotations, MilLung combines nodule-level supervision with multiple instance learning, while AutoLung learns representations with an autoencoder and requires only patient-level malignancy labels. The evaluation uses LIDC-IDRI under two malignancy thresholds. At a threshold greater than three, the best MilLung variant matches DeepLung's mean area under the receiver operating characteristic curve of 0.83, while AutoLung reaches 0.80. AutoLung performs less well at the lower threshold, showing that the effect of reduced supervision depends on the task. The study clarifies the trade-off between annotation requirements and predictive performance when selecting a patient-level screening approach.
    \end{quote}
\end{itemize}

\section{Patents}

\begin{itemize}
    \item \textbf{European Patent application EP24461637.1} ``A computer implemented method for identifying a microorganism in a blood and a data processing system therefor''. European Patent Office, 2024, patent pending
\end{itemize}

\section{Ongoing research}

\begin{itemize}
    \item \textbf{Reinforcement learning for visual exploration in real-world environments.} A. Pardyl, J. Mikołajczyk, J. Matusiak, H. Zientata, D. Gaweł, I. Alwasiak, B. Zieliński, M. Cygan, M. Wołczyk
    
    Summary:
    \begin{quote}
    Training agents for visual exploration requires diverse navigation experience and a way to transfer learned behaviour to physical environments. This project extends the FlySearch~\pubcite{pardyl2025flysearch} benchmark towards reinforcement learning for vision-language model-based agents operating in the real world. The planned approach uses the GRPO algorithm~\cite{shao2024deepseekmath} to fine-tune models with 4--11 billion parameters on FlySearch-RealTime, a simplified version of the benchmark, and Gaussian splat scenes reconstructed from real-world flight data. Procedurally generated tasks will provide varied navigation experience, while the reconstructed scenes are intended to reduce the visual gap between simulation and deployment. Evaluation will use the original FlySearch benchmark and a physical UAV in a real-world setting. The intended contribution is a training framework for learning exploration policies that transfer from simulated tasks to physical navigation.
    \end{quote}

    \item \textbf{Self-supervised concept discovery.} A. Pardyl, S. Gairola, S. Rao, B. Schiele, B. Zieliński, D.~\mbox{Rymarczyk}
    
    Summary:
    \begin{quote}
    Self-supervised visual representations support a range of downstream tasks, but their features can be difficult to interpret. This project aims to learn discrete, spatially grounded concepts corresponding to visual features such as object parts or texture patterns without manual concept annotations. Inspired by PDiscoFormer~\cite{aniraj2024pdiscoformer} and CFM~\cite{wittenmayer2026cfm}, the proposed approach extends the DINOv2 family~\cite{oquab2024dinov2} with a spatial concept learning pathway and a concept attribute quantisation stack to represent both general object characteristics and their variations. The resulting decoded concept embedding is intended to replace the \textsc{cls} token of a vision transformer in downstream tasks. The project will examine whether the learned concepts capture coherent visual features while retaining useful information for downstream tasks. Its intended contribution is an interpretable vision foundation backbone that connects image-level representations to discrete concepts and their spatial locations.
    \end{quote}
    
\end{itemize}

\section{International research internships}
\begin{itemize}
    \item \textbf{Max Planck Institute for Informatics}, June--September 2026, Computer Vision Group, supervised by Prof. Dr. Bernt Schiele
\end{itemize}

\section{Scientific grants}
% Grant numbers are long and unbreakable; \sloppy avoids overfull lines here.
\begingroup
\sloppy
\subsection{As a principal investigator}
\begin{itemize}
    \item \textbf{National Science Centre Preludium} ``Where to look next -- guiding active visual exploration with
    internal model uncertainty'', 2024--2025, grant number: 2023/49/N/ST6/02465
    \item \textbf{Foundation for Polish Science START scholarship}, 2026--2027, grant number: 59.2026
\end{itemize}

\subsection{As a participant / co-investigator}
\begin{itemize}
    \item \textbf{National Science Centre Sonata Bis} ``Sustainable computer vision for autonomous machines'', 2024--2027, grant number: 2023/50/E/ST6/00469 (principal investigator: Bartosz Zieliński)
    \item \textbf{Foundation for Polish Science Proof of Concept} ``Intelligent Fight Against Sepsis - Innovative Use of Artificial Intelligence in Rapid Microbiological Diagnosis of Sepsis'', 2024--2025, grant number: FENG.02.07-IP.05-0068/23 (principal investigator: Monika Brzychczy-Włoch)
    \item \textbf{European Commission Horizon Europe} ``European Lighthouse of AI for Sustainability (ELIAS)'', 2024--2027, grant number: 101120237 (principal investigator: Massimiliano Pontil)
    \item \textbf{Foundation for Polish Science Team-Net} ``Bio-inspired artificial neural networks'', 2022--2023, grant number: POIR.04.04.00-00-14DE/18-00 (principal investigator: Jacek Tabor)
    \item \textbf{National Centre for Research and Development Fast Track} ``X-rAI: Diagnostic viewer for computer-assisted radiology using Artificial Intelligence'', 2021--2023, grant number: POIR.01.01.01-\allowbreak 00-\allowbreak 1666/20 (principal investigator: Zbisław Tabor)
\end{itemize}

\endgroup

\section{Other activities}

\subsection{Acting as a reviewer}
\begin{itemize}
    \item {Annual Conference on Neural Information Processing Systems (NeurIPS) 2026}
    \item {International Journal of Computer Vision (IJCV) 2026}
    \item {IEEE Transactions on Pattern Analysis and Machine Intelligence (TPAMI) 2026}
    \item {IEEE Robotics and Automation Letters (RA-L) 2026}
    \item {CVPR 2025 Workshop on Computer Vision for Drug Discovery (CVDD)}
    \item {IEEE/CVF Winter Conference on Applications of Computer Vision (WACV) 2025}
    \item {ACM International Conference on Knowledge Discovery and Data Mining (KDD) 2024}
    \item {IEEE/CVF Winter Conference on Applications of Computer Vision (WACV) 2024}
\end{itemize}

\subsection{Doctoral schools and scientific events}
\subsubsection{As an organiser}
\begin{itemize}
    \item Machine Learning Summer School on Reliability \& Safety (MLSS$^{R\&S}$) 2026, Kraków, Poland
    \item ELLIS Doctoral Symposium (EDS) on Robust AI 2025, Warsaw, Poland
    \item {\emergencystretch=2em Machine Learning Summer School on Drug and Materials Discovery (MLSS$^{D}$) 2025, Kraków,~Poland\par}
    \item Machine Learning Summer School on Applications in Science (MLSS$^{S}$) 2023, Kraków, Poland
\end{itemize}

\subsubsection{As a presenter}
\begin{itemize}
    \item {\emergencystretch=2em ML in PL Conference 2025, Warsaw, Poland, ``FlySearch: Exploring how vision-language \mbox{models explore}''\par}
    \item ML in PL Conference 2024, Warsaw, Poland, ``AdaGlimpse: Active Visual Exploration with Arbitrary Glimpse Position and Scale''
    \item SFI IT Academic Festival 2024, Kraków, Poland, ``AI for autonomous agents''
\end{itemize}

\subsubsection{As a participant}
\begin{itemize}
    \item Mediterranean Machine Learning (M2L) Summer School 2025, Split, Croatia
    \item International Computer Vision Summer School (ICVSS) 2025, Scicli, Italy
    \item ELIAS-ELLIS-VISMAC Winter School 2025, Brunico, Italy
    \item ELLIS Doctoral Symposium (EDS) on AI \& Sustainability 2024, Paris, France
    \item Mediterranean Machine Learning (M2L) Summer School 2023, Thessaloniki, Greece
\end{itemize}

\subsection{Acting as a PhD Student Representative}
\begin{itemize}
    \item Member of the Audit Committee of the PhD Student Association of the Jagiellonian University, 2025--2026
    \item Member of the PhD Student Council of the Doctoral School of Exact and Natural Sciences of the Jagiellonian University, 2024--2026
    \item PhD Student Representative to the Standing Rector's Committee on Statutory and Legal Affairs, 2022--2026
    \item PhD Student Representative to the Council of the Faculty of Mathematics and Computer Science of the Jagiellonian University, 2025--2026
\end{itemize}
